\documentclass[12pt]{article}
% \usepackage[top=1in,left=1in, right = 1in, footskip=1in]{geometry}
\usepackage[top=1in,footskip=1in]{geometry}

\usepackage{graphicx}
\usepackage{xspace}
%\usepackage{adjustbox}

\usepackage{pdflscape}

\usepackage{grffile}

\newcommand{\comment}{\showcomment}
%% \newcommand{\comment}{\nocomment}

\newcommand{\showcomment}[3]{\textcolor{#1}{\textbf{[#2: }\textsl{#3}\textbf{]}}}
\newcommand{\nocomment}[3]{}

\newcommand{\jd}[1]{\comment{cyan}{JD}{#1}}
\newcommand{\swp}[1]{\comment{magenta}{SWP}{#1}}
\newcommand{\bmb}[1]{\comment{blue}{BMB}{#1}}
\newcommand{\djde}[1]{\comment{red}{DJDE}{#1}}

\newcommand{\eref}[1]{Eq.~(\ref{eq:#1})}
\newcommand{\fref}[1]{Fig.~\ref{fig:#1}}
\newcommand{\Fref}[1]{Fig.~\ref{fig:#1}}
\newcommand{\sref}[1]{Sec.~\ref{#1}}
\newcommand{\frange}[2]{Fig.~\ref{fig:#1}--\ref{fig:#2}}
\newcommand{\tref}[1]{Table~\ref{tab:#1}}
\newcommand{\tlab}[1]{\label{tab:#1}}
\newcommand{\seminar}{SE\mbox{$^m$}I\mbox{$^n$}R}

\usepackage{amsthm}
\usepackage{amsmath}
\usepackage{amssymb}
\usepackage{amsfonts}
\usepackage[utf8]{inputenc} % make sure fancy dashes etc. don't get dropped

\usepackage{lineno}
\linenumbers

\usepackage[pdfencoding=auto, psdextra]{hyperref}

\usepackage{natbib}
\bibliographystyle{chicago}
\date{\today}

\usepackage{xspace}
\newcommand*{\ie}{i.e.\@\xspace}

\usepackage{color}

\newcommand{\Rx}[1]{\ensuremath{{\mathcal R}_{#1}}\xspace} 
\newcommand{\RR}{\ensuremath{{\mathcal R}}\xspace}
\newcommand{\Rres}{\Rx{\mathrm{res}}}
\newcommand{\Rinv}{\Rx{\mathrm{inv}}}
\newcommand{\Rhat}{\ensuremath{{\hat\RR}}}
\newcommand{\Rt}{\ensuremath{{\mathcal R}(t)}\xspace}
\newcommand{\tsub}[2]{#1_{{\textrm{\tiny #2}}}}
\newcommand{\dd}[1]{\ensuremath{\, \mathrm{d}#1}}
\newcommand{\dtau}{\dd{\tau}}
\newcommand{\dx}{\dd{x}}
\newcommand{\dsigma}{\dd{\sigma}}

\newcommand{\rx}[1]{\ensuremath{{r}_{#1}}\xspace} 
\newcommand{\rres}{\rx{\mathrm{res}}}
\newcommand{\rinv}{\rx{\mathrm{inv}}}

\newcommand{\psymp}{\ensuremath{p}} %% primary symptom time
\newcommand{\ssymp}{\ensuremath{s}} %% secondary symptom time
\newcommand{\pinf}{\ensuremath{\alpha_1}} %% primary infection time
\newcommand{\sinf}{\ensuremath{\alpha_2}} %% secondary infection time

\newcommand{\psize}{{\mathcal P}} %% primary cohort size
\newcommand{\ssize}{{\mathcal S}} %% secondary cohort size

\newcommand{\gtime}{\tau_{\rm g}} %% generation interval
\newcommand{\gdist}{g} %% generation-interval distribution
\newcommand{\idist}{\ell} %% incubation-period distribution

\newcommand{\total}{{\mathcal T}} %% total number of serial intervals

\usepackage{lettrine}

\newcommand{\dropcapfont}{\fontfamily{lmss}\bfseries\fontsize{26pt}{28pt}\selectfont}
\newcommand{\dropcap}[1]{\lettrine[lines=2,lraise=0.05,findent=0.1em, nindent=0em]{{\dropcapfont{#1}}}{}}

\begin{document}

\begin{flushleft}{
	\Large
	\textbf\newline{
	  Strategic coexistence theory for evolutionary games, another case study
	}
}
\newline
\\
Sang Woo Park\textsuperscript{1,2,*}
\\
\bigskip
\textbf{1} School of Biological Sciences, Seoul National University, Seoul, Korea\\
\textbf{2} Institute for Data Innovation in Science, Seoul National University, Seoul, Korea
*Corresponding author: sangwoopark@snu.ac.kr
\end{flushleft}

So far, we have focused on simple deterministic 2×2 games. 
However, many ecological and evolutionary processes are intrinsically stochastic.
Recently, \cite{wang2025evolution} showed that environmental fluctuations can promote coexistence in games that would otherwise lead to competitive exclusion in a deterministic environment. 
From the perspective of SCT, this suggests that stochasticity can stabilize competition by promoting the recovery of strategies when rare. 
Interestingly, \cite{wang2025evolution} also showed that stochasticity can generate bistable coexistence and multiple internal equilibria.
Thus, although stochasticity can stabilize competition near the boundaries, it can also generate destabilizing frequency dependence in the interior. 
Here, we use SCT to separate these effects and clarify how stochasticity can both stabilize and destabilize strategic competition.

Following \cite{wang2025evolution}, we consider a $2×2$ game in the presence of environmental noise. 
The stochastic payoff matrix is given by
\begin{equation}
\begin{pmatrix}
R & S \\
T & P
\end{pmatrix} +
\frac{\zeta}{\sqrt{s}} \begin{pmatrix}
\sigma_R & \sigma_S \\
\sigma_T & \sigma_P
\end{pmatrix}, 
\end{equation}
where the first matrix represents deterministic payoffs,
$\zeta$ represents a noise term with zero mean and unit variance,
and the second matrix determines the intensity of environmental stochasticity.
Assuming weak selection $s<<1$, the evolution of two competing strategies in a large population can be described by the following equation:
\begin{equation}
\frac{dx}{dt} = s x (1-x) \left(\pi_C - \pi_D + \left(\frac{1}{2} - x \right) (\sigma_C - \sigma_D)^2 \right),
\end{equation}
where $\pi_C = Rx + S(1-x)$ and $\pi_D = T x + P(1-x)$ and $\sigma_C = \sigma_R x + \sigma_S (1-x)$ and $\sigma_D = \sigma_T x + \sigma_P (1-x)$.
For this system, the invasion growth rates of each strategy when rare can be written as:
\begin{align}
r_C &= S - P + \frac{1}{2}(\sigma_S - \sigma_P)^2,\\
r_D &= T - R + \frac{1}{2}(\sigma_T - \sigma_R)^2. 
\end{align}
This shows that environmental stochasticity always increases the invasion growth rates, meaning that it is is intrinsically stabilizing.

Then, an extended SCT decompsition shows that niche overlap and fitness difference can be written as:
\begin{align}
\rho &= \exp\left(\frac{R+P-S-T}{2} - \frac{(\sigma_R - \sigma_T)^2 + (\sigma_S - \sigma_P)^2}{4} \right)\\
\frac{\kappa_D}{\kappa_C} &= \exp\left(\frac{T+P-R-S}{2} + \frac{(\sigma_R - \sigma_T)^2 - (\sigma_S - \sigma_P)^2}{4} \right)
\end{align}
Writing $\hat{\rho}$ and $\hat{\kappa}_D/\hat{\kappa}_C$ to denote niche overlap and fitness difference in the absence of stochasticity and takign logs of both sides for convenience, we have:
\begin{align}
\Delta \log \rho = \log \rho -  \log  \hat{\rho}&=- \frac{(\sigma_R - \sigma_T)^2 + (\sigma_S - \sigma_P)^2}{4}\\
\Delta \log \frac{\kappa_D}{\kappa_C} = \log \frac{\kappa_D}{\kappa_C}-\log \frac{\hat{\kappa}_D}{\hat{\kappa}_C} &= \frac{(\sigma_R - \sigma_T)^2 - (\sigma_S - \sigma_P)^2}{4}.
\end{align}
Therefore, stochasticity always reduces niche overlap and therefore increases niche difference.
In contrast, stochasticity can either amplify or reduce fitness differences depending on which strategy’s invasion rate is more strongly affected by environmental fluctuations.

This immediately raises a question of whether stochasticity can push the system out of the coexistence regime.
However, this is not possible.


\subsection*{Supp}

Consider a competition of two strategies under stochastic payoffs:
\begin{equation}
\begin{pmatrix}
R & S \\
T & P
\end{pmatrix} +
\frac{\zeta}{\sqrt{s}} \begin{pmatrix}
\sigma_R & \sigma_S \\
\sigma_T & \sigma_P
\end{pmatrix}, 
\end{equation}
Assuming weak selection $s<<1$, the evolution of two competing strategies in a large population can be described by the following equation:
\begin{equation}
\frac{dx}{dt} = s x (1-x) \left(\pi_C - \pi_D + \left(\frac{1}{2} - x \right) (\sigma_C - \sigma_D)^2 \right),
\end{equation}
To simplify, we write
\begin{equation}
\frac{dx}{dt} = s x (1-x) G(x),
\end{equation}
where
\begin{equation}
G(x) = \left(\pi_C - \pi_D + \left(\frac{1}{2} - x \right) (\sigma_C - \sigma_D)^2 \right).
\end{equation}
Thus, the sign of $G(x)$ at boundary conditions determines whether a strategy can invade when the competing strategy is at equilibrum.

Note that we cannot immediately apply the original SCT presented in the main text here because $G(x)$ does not uniquely separate into two payoff measures.
Instead, we can still utilize the invasion growth rate when rare and follow the same procedure.
In particular, the mutual invasion condition at monomorphic boundary can be written as:
\begin{align}
0 &< G(0)\\
0 &< -G(1)
\end{align}
Exponentiating both, we can combine them into a single inequality:
\begin{align}
\exp(G(1)) < 1< \exp(G(0))
\end{align}
Dividing both sides by $\exp((G(1) + G(0))/2)$, we get:
\begin{align}
\exp\left(\frac{G(1)-G(0)}{2}\right)<  \exp\left(-\frac{G(1)+G(0)}{2}\right) < \exp\left(\frac{G(0)-G(1)}{2}\right),
\end{align}
which gives us expressions for niche overlap and fitness difference:
\begin{align}
\rho &= \exp\left(\frac{G(1)-G(0)}{2}\right)\\
\frac{\kappa_D}{\kappa_C} &= \exp\left(-\frac{G(1)+G(0)}{2}\right)
\end{align}
Then, by substituting expressions for $G(0)$ amd $G(1)$, we obtain following expressions for niche and fitness difference:
\begin{align}
1- \rho &= 1-\exp\left(\frac{R+P-S-T}{2} - \frac{(\sigma_R - \sigma_T)^2 + (\sigma_S - \sigma_P)^2}{4} \right)\\
\frac{\kappa_D}{\kappa_C} &= \exp\left(\frac{T+P-R-S}{2} + \frac{(\sigma_R - \sigma_T)^2 - (\sigma_S - \sigma_P)^2}{4} \right)
\end{align}

\bibliography{coexistence-strategy}

\end{document}
